\newglossaryentry{term1}{
    name={Нагрузочное тестирование},
    description={вид тестирования производительности, направленный на определение поведения системы под ожидаемой и повышенной нагрузкой}
}
\newglossaryentry{term2}{
    name={Система управления базами данных},
    description={программное обеспечение для создания, ведения и управления базами данных, обеспечения доступа к ним и контроля целостности}
}
\newglossaryentry{term3}{
    name={Пропускная способность},
    description={количество успешно обработанных транзакций в единицу времени, измеряемое в TPS (Transactions Per Second)}
}
\newglossaryentry{term4}{
    name={Задержка},
    description={время от отправки запроса к базе данных до получения ответа, измеряемое в миллисекундах}
}
\newglossaryentry{term5}{
    name={Процентиль},
    description={значение, ниже которого находится определённый процент измерений в наборе данных (например, p95 означает, что 95\% запросов выполнены быстрее данного значения)}
}
\newglossaryentry{term6}{
    name={Виртуальный пользователь},
    description={эмулированная сущность, выполняющая транзакции в ходе нагрузочного теста для имитации реальной рабочей нагрузки}
}
\newglossaryentry{term7}{
    name={Сценарий тестирования},
    description={набор параметризованных SQL-запросов с заданным соотношением операций, моделирующих рабочую нагрузку на СУБД}
}
\newglossaryentry{term8}{
    name={Временной ряд},
    description={последовательность значений метрик производительности, зафиксированных через равные интервалы времени в ходе выполнения теста}
}
\newglossaryentry{term9}{
    name={Снапшот состояния},
    description={фиксация состояния базы данных (количество строк в таблицах) перед выполнением теста для последующей верификации целостности}
}
\newglossaryentry{term10}{
    name={Стресс-тестирование},
    description={вид тестирования производительности для определения границ работоспособности системы при экстремальных нагрузках, отличающийся от нагрузочного тестирования}
}
\newglossaryentry{term11}{
    name={Конфигурация сценария нагрузки},
    description={конкретный набор параметризованных SQL-запросов, связанный с профилем схемы и шаблоном сценария; определяет состав, веса и параметры запросов при выполнении тестового прогона}
}
\newglossaryentry{term12}{
    name={Логическая база данных},
    description={именованная группа физических подключений к СУБД, объединённых единым профилем схемы; используется как единица конфигурации при запуске нагрузочного тестирования нескольких экземпляров СУБД}
}
\newglossaryentry{term13}{
    name={Профиль схемы},
    description={описание структуры данных, общей для группы совместимых баз данных одной предметной области; служит основой для автоматической генерации конфигураций сценариев нагрузки}
}
\newglossaryentry{term14}{
    name={Размер эффекта},
    description={стандартизированная мера практической значимости различий между двумя выборками; в настоящей работе применяется коэффициент Коэна~$d$, выражающий величину различия в единицах стандартного отклонения}
}
